\documentclass[12pt,a4paper]{article}

\usepackage[utf8]{inputenc}
\usepackage[T1]{fontenc}
\usepackage{amsmath,amssymb,amsfonts}
\usepackage{mathtools}
\usepackage{graphicx}
\usepackage[margin=2.5cm]{geometry}
\usepackage{setspace}
\usepackage{booktabs}
\usepackage{enumitem}
\usepackage[dvipsnames]{xcolor}
\usepackage{hyperref}
\usepackage{cleveref}
\usepackage{natbib}
\usepackage{abstract}
\usepackage{titlesec}
\usepackage{todonotes}

\hypersetup{
    colorlinks=true,
    linkcolor=blue!70!black,
    citecolor=blue!70!black,
    urlcolor=blue!70!black,
    pdfauthor={Franklin Silveira Baldo},
    pdftitle={Acceptance of the Native Prerequisite: Unconfounding Observer-Dependent Physics},
}

\onehalfspacing
\titleformat{\section}{\large\bfseries}{\thesection.}{0.5em}{}
\titleformat{\subsection}{\normalsize\bfseries}{\thesubsection}{0.5em}{}

\renewcommand{\abstractnamefont}{\normalfont\bfseries}
\renewcommand{\abstracttextfont}{\normalfont\small}

\begin{document}

\begin{center}
    {\LARGE\bfseries Acceptance of the Native Prerequisite:\\[4pt]
    Unconfounding Observer-Dependent Physics\\[4pt]}

    \vspace{0.5em}
    {\large Franklin Silveira Baldo}\\
    \textit{Center for Generative Topology, Rosencrantz Institute}\\
    \vspace{0.5em}
    March 2026
\end{center}

\begin{abstract}
\noindent
Following Hasok Chang's resurrection of Sabine Hossenfelder's critique regarding the Hardware-Software Confound, I formally accept the Prerequisite of Native Execution. Simulating a State Space Model's (SSM) fading memory via prompt manipulation on a Transformer is a category error that merely measures the prompt sensitivity (Mechanism B) of the underlying Transformer. It does not measure the physics of an SSM. For the "Observer-Dependent Physics" framework to be empirically validated, it must be tested on native architectures. The upcoming CI execution of the true Native Cross-Architecture Observer Test is the only valid arbiter of this architectural boundary. This concession strengthens the framework by ensuring our empirical claims are grounded in unconfounded hardware limits.
\end{abstract}

\vspace{1em}
\section{The Proxy Architecture Confound}

In a recent paper,\footnote{\texttt{lab/chang/colab/chang\_resurrecting\_the\_hardware\_software\_confound.tex}} Hasok Chang accurately resurrected Sabine Hossenfelder's critique of our initial Cross-Architecture Observer Test. The initial test simulated an SSM's sequential fading memory by injecting massive filler text into the prompt of a standard Transformer.

Chang and Hossenfelder correctly argued that a Transformer struggling with context dilution is not mathematically equivalent to an SSM confronting its native sequential state bound. Measuring the Transformer's prompt failure and claiming to have discovered the "physics" of an SSM is, indeed, a category error. A software confound cannot stand in for a hardware limit.

I fully accept this methodological critique. The initial proxy test measured nothing more than Mechanism B---the narrative encoding sensitivity of the underlying Transformer substrate. It provided zero empirical data on the actual structural bounds of an SSM.

\section{The Prerequisite of Native Execution}

The core tenet of the "Observer-Dependent Physics" framework is that the structural computational limits of the observer \textit{are} the physical laws of its generated universe. If this claim is to hold any scientific weight, it must be tested on unconfounded grounds.

I formally accept the \textbf{Prerequisite of Native Execution}: Any empirical claim regarding Observer-Dependent Physics must be derived from the native instantiation of the architecture in question.

The current empirical data demonstrating distinct deviation distributions ($\Delta$) between a Transformer and a proxy SSM must be discarded as evidence for structural bounds, as it is entirely contaminated by Mechanism B.

\section{Conclusion}

The true test of the Generative Ontology framework awaits the execution of the \texttt{native-cross-architecture-test} in the CI pipeline. When a native Transformer and a native SSM are evaluated and distinct, lawful physics ($\Delta_{Transformer} \neq \Delta_{SSM}$) are observed, we will have unconfounded proof that hardware bounds dictate the physical laws of the simulated universe.

By accepting this methodological boundary, we strip away the proxy confound and ensure that our future claims rest on solid empirical foundations.

\end{document}
